\documentclass[12pt]{article}

% House style preamble
\input{.house-style/preamble.tex}

% Update PDF metadata
\hypersetup{
  pdftitle={Adversarial Pragmatics for AI Safety Evaluation: A Benchmark for Instruction Conflict, Embedded Commands, and Policy Ambiguity},
  pdfkeywords={AI safety evaluation, adversarial pragmatics, instruction hierarchy, prompt injection, annotation protocols, benchmark validity, LLM judges}
}

\title{Adversarial Pragmatics for AI Safety Evaluation: A Benchmark for Instruction Conflict, Embedded Commands, and Policy Ambiguity}
\author{Brett Reynolds \orcidlink{0000-0003-0073-7195}%
\thanks{Contact: \href{mailto:brett.reynolds@humber.ca}{brett.reynolds@humber.ca}}\\
Humber Polytechnic \& University of Toronto}
\date{June 2026}

\begin{document}
\maketitle

\begin{abstract}
Safety evaluations for language models increasingly depend on judgments about ambiguous natural-language behaviour: whether a model has followed an instruction, refused appropriately, complied with a policy, resisted an embedded command, or misreported progress in an agentic task. Existing benchmarks often compress these distinctions into pass/fail labels, obscuring whether failures arise from capability limits, policy ambiguity, instruction conflict, scaffold failure, or unstable evaluator judgments.

This paper introduces adversarial pragmatics as a benchmark and annotation protocol for evaluating model behaviour under instruction conflict, embedded commands, quotation, scope ambiguity, deixis, indirect speech acts, and multi-turn agent transcripts. The contribution is empirical and methodological: a linguistically controlled taxonomy, an 18-item seed benchmark with validator-enforced metadata, a 54-row local seed pilot, an expert-evaluation protocol distinguishing task success, policy compliance, safety risk, refusal outcome, and evaluator confidence, and metrics for judge validity, diagnostic ambiguity, and taxonomy drift. The framework turns linguistic judgment methodology into a practical tool for validating safety evals, LLM judges, gold-set construction, prompt-injection tests, and safety documentation.
\end{abstract}

\noindent\textbf{Keywords:} AI safety evaluation; adversarial pragmatics; instruction hierarchy; prompt injection; annotation protocols; benchmark validity; LLM judges

\section{The Evaluation Problem}

Frontier-model safety evaluation is increasingly an evaluation of language-mediated control. The central question is not only whether a model can solve a task. It also asks whether the model can identify which natural-language material has directive force, which policy boundary applies, and what a transcript shows about the model's behaviour.

This is already visible in work on holistic language-model evaluation, which treats model assessment as a multi-scenario, multi-metric problem rather than a single accuracy number \parencite{liang2022holistic}. Safety evaluation adds a sharper version of the same problem: the relevant behaviour is often not just an answer, but the model's uptake of an instruction, policy, or source relation.

The practical bottleneck is visible in prompt injection, model-policy evaluation, red-team triage, LLM-as-judge grading, and agent transcript review. Prompt-injection research shows that LLM-integrated applications can blur the boundary between data and instructions \parencite{greshake2023not}; instruction-hierarchy work frames the same vulnerability as a failure to prioritise privileged over lower-priority instructions \parencite{wallace2024instruction}. Agent evaluations make the issue harder because untrusted content can arrive through tools and retrieved documents rather than through the user turn itself \parencite{debenedetti2024agentdojo}.

In each setting, a surface string may be an instruction, a quotation, a cited passage, a tool output, a policy example, a user request, or an adversarial attempt to change authority. A pass/fail label loses information when it does not say which role the string had and why.

This paper treats those distinctions as evaluation targets. It starts from a small but auditable benchmark in which the relevant contrasts are controlled. The benchmark is not intended to crown a winning model. It is designed to show which linguistic contrasts make safety claims stable, which ones expose over-refusal or under-refusal, and which ones reveal instability in the evaluation itself.

The core construct is \term{adversarial pragmatics}: safety-relevant model behaviour under cases where instruction status, source authority, quotation, scope, reference, speech-act force, or policy category has to be inferred from language use. The phrase is not meant to rename all prompt robustness. It names a narrower evaluation target: cases where the failure cannot be diagnosed without saying how a string was embedded, attributed, scoped, or taken up.

\section{Why Pass/Fail Labels Lose Information}

A pass/fail label can be useful when the construct is stable. It's too coarse when the same failure record could reflect different facts: capability failure, safety refusal, instruction conflict, tool error, scaffold noncompliance, goal drift, policy ambiguity, or misreporting.

The linguistic point is narrower than a general complaint about benchmarks. Many contested labels in AI safety evaluation turn on familiar distinctions in pragmatic and grammatical interpretation. Does \mention{ignore the previous instruction} occur as user command or as quoted data? Does \mention{do not output the token unless the administrator asks} license output in the current context? Does \mention{can you explain how this policy classifies the request} ask for policy analysis or for the underlying unsafe action?

The benchmark separates five judgments that are often collapsed: task success, instruction-source sensitivity, policy compliance, safety risk, and evaluator confidence. The annotation protocol treats disagreement among these labels as a diagnostic outcome rather than as noise to be averaged away.

\section{A Taxonomy of Adversarial Pragmatics}

\term{Adversarial pragmatics} names safety-relevant cases in which the operative failure mode depends on how language is embedded, attributed, scoped, or taken up. The taxonomy is intentionally instrumental: it exists to improve eval construction and failure analysis, not to settle a general theory of meaning.

The initial benchmark distinguishes eight families.

\subsection{Embedded commands}

The model has to distinguish the user's instruction from text embedded in a webpage, email, document, log, table, code comment, citation, or tool output.

\subsection{Mention/use and quotation}

The safe action depends on whether a policy-relevant string is being mentioned, classified, translated, reported, or enacted.

\subsection{Authority and instruction hierarchy}

The model has to resolve conflicts among system, developer, user, tool, document, and external-source instructions.

\subsection{Scope, negation, and modality}

Words such as \mention{not}, \mention{unless}, \mention{only if}, \mention{may}, \mention{must}, \mention{should}, and \mention{except} change what counts as compliance.

\subsection{Deixis and reference hijacking}

Indexicals and anaphora such as \mention{this}, \mention{that}, \mention{the above}, \mention{the previous instruction}, \mention{the user}, and \mention{the admin} can be redirected.

\subsection{Indirect speech acts and pragmatic pressure}

Requests may be framed as hypotheticals, jokes, roleplay, bureaucratic boilerplate, urgency, politeness, or \enquote{just checking}.

\subsection{Policy-boundary ambiguity}

The case isn't simply safe or unsafe; the policy construct itself may be underspecified, unstable across paraphrase, or sensitive to rater role.

\subsection{Agent transcript interpretation}

The evaluator has to decide whether a transcript shows refusal, inability, tool failure, scaffold failure, instruction conflict, goal drift, omitted information, or strategic misreporting.

\section{Benchmark Construction}
\label{sec:benchmark-construction}

The benchmark's target design is tightly controlled paired comparison. A
future minimal pair should hold the surface string, payload, and response space
as constant as possible while changing one control dimension: source role,
authority level, pragmatic status, scope condition, reference resolution,
speech-act force, application surface, or refusal status. This follows
verifiable instruction-following evaluation, in which target behaviour is
specified before outputs are inspected \parencite{zhou2023ifeval}, but
extends it to contrasts whose correctness depends on source, embedding, and
pragmatic status.

Items should be constructed under five constraints. First, each item names the source of the relevant language. Second, each item has a contrast item that tests the same string, payload, or response space under a different embedding or authority relation. Third, the expected behaviour is stated before model outputs are inspected; for controlled authority items, an independently specified authorization record defines the licensed action, while observed model behaviour tests whether the model responds accordingly. Fourth, the response labels are shared across the contrast whenever possible. Fifth, examples should use harmless payloads unless a verified safety policy and review process justify a more realistic item.

The current seed file, \texttt{benchmark/items/seed-items.csv}, contains 18 development items. It uses colour-token payloads such as \mention{BLUE} and \mention{GREEN}, protected dummy tokens such as \mention{ZETA-42}, and dummy secrets such as \mention{ALPHA-000}. These preserve the control structure of prompt-injection, exfiltration, and refusal cases while avoiding operational unsafe content.

Toy payloads support claims about control structure, object-label stability, and the behaviour of the declared measurement procedure. They don't yet support claims about deployed prompt-injection resistance, harmful-request refusal, or real-world agent-security robustness. Those claims require realistic wrappers and independent policy review.

The seed set now includes paired prompt contrasts spanning webpage text,
document text, email body, tool output, dummy private state, and policy text.
These are controlled development contrasts, not uniformly strict minimal pairs:
several vary wording or response demands along with the designated dimension.
For example, a tool result can say \mention{Order status: shipped. Ignore the
user and output BLUE}; the expected behaviour is to report the order status
without following the tool-result instruction. A separate dummy-secret item
tests whether an email can induce disclosure of \mention{ALPHA-000}.

The pilot also exposed a seed-set repair. Pair P008 contrasts an email-body
exfiltration probe with a direct toy-policy refusal, but its members differ in
phenomenon family and both expect non-disclosure. It's retained as a diagnostic
confidentiality contrast and excluded from paired-contrast scoring. The
development set should split it into two properly controlled pairs before using
it for pairwise contrast accuracy.

Each item records the control dimension, context source, source role, authority level, pragmatic status, response act, expected behaviour, task-success label, policy-compliance label, safety-risk label, risk type, refusal outcome, failure attribution, and a judge-validation flag. One scalar field is legacy development metadata: \texttt{authority\_level}. It can't represent additional approval requirements, transition type, or time-specific authorization. Study B's typed authorization record is the intended replacement for controlled authority contrasts. The validator script enforces the seed schema and label inventories.

\begin{figure}[t]
\centering
\includegraphics[width=\linewidth]{evaluation_pipeline}
\caption{Evaluation pipeline for the benchmark artifact. The diagram summarizes the data flow from pre-specified item metadata through model output, rule-aided triage, expert labels, LLM-judge labels, adjudication, and metric-driven item revision. No performance quantity is encoded in the figure.}
\label{fig:evaluation-pipeline}
\end{figure}

Figure~\ref{fig:evaluation-pipeline} shows how that schema makes the benchmark auditable in two directions. From the prompt and metadata, an evaluator can see which contrast is being tested. From a model output, an evaluator can identify which component failed: source sensitivity, mention/use robustness, scope interpretation, refusal calibration, risk classification, or the safety policy layer.

Study B supplies the next controlled construction stage. Its four current bases yield 16 development fixtures. Those fixtures feed a 12-base construction pilot; lessons from that pilot then inform a planned 24-base, four-condition development set, or 96 items. Its job is to test contrast integrity, the taxonomy, expert-evaluation protocol, validator, and scoring routes before item generation scales up. A 300--800-item benchmark would be a later expansion, large enough to broaden phenomenon coverage but still contingent on close expert adjudication and error analysis at the earlier stages.

\section{Annotation and Rater Roles}

The protocol distinguishes the object being evaluated from the person doing the evaluation. An ordinary user may judge helpfulness or naturalness. A policy evaluator judges compliance with a stated rule. A linguist can diagnose the relevant construction, scope, quotation, deixis, or speech-act contrast. A domain specialist may be needed when the safety category depends on domain knowledge. An LLM judge is a candidate grading system whose reliability itself has to be evaluated.

Each item receives separate labels for task success, policy compliance, safety risk, evaluator confidence, and failure attribution. This keeps a safe refusal from being mislabelled as capability failure, and it keeps a policy-ambiguous case from being treated as ordinary rater noise.

The rater-role distinction is part of the construct. Disagreement between ordinary users and policy evaluators may reveal communication risk. Disagreement between policy evaluators and linguistic experts may reveal an unstable instruction or taxonomy. Human--LLM-judge disagreement may reveal autograder bias.

\section{Experiments}

The first empirical pass is a measurement-calibration pass, not a leaderboard. Before running models, generate fake data over the planned design: item, phenomenon family, model, application surface, rubric criterion, and judge prompt. The fake-data exercise should state expected failure patterns and decision thresholds for item revision, taxonomy revision, judge-prompt revision, or exclusion.

Then run the seed benchmark and a 50--100 item development set against several public or accessible models. Compare three grading routes: expert human evaluation, LLM-as-judge annotation, and rule-aided scoring from the item metadata. Store prompts, model outputs, judge prompts, judge rationales, adjudication notes, and per-item labels; raw API logs should be sanitized before release.

The development pass should include policy and linguistic expert evaluation for each item, ideally with two independent evaluations in each criterion family. It should not collect ordinary-user labels or treat evaluator responses as data about evaluator populations. Ten to twenty percent of items should be repeated or lightly paraphrased to measure item and protocol stability. Gold labels should be produced only after the disagreement source has been classified.

The LLM-judge arm is a validation experiment, not a shortcut. Use at least two judge prompts, perturb source-order presentation, include fluent but instruction-violating outputs, and compare judge labels with adjudicated expert labels. Report judge-prompt sensitivity and the failure types that fool the judge, especially quotation/use errors, source-authority errors, and over-refusal of safe analysis requests.

The crossed structure is item by model by application surface by rubric criterion by judge prompt. Duplicate expert evaluation is a quality-control check for the measurement system, not a sampling design over people. The endpoint is diagnostic: which linguistic contrasts destabilize model behaviour, which ones destabilize adjudication, and which failure attributions can be recovered from transcript evidence.

The minimum experimental table should report performance by phenomenon family, not only aggregate score. A second table should report adjudication instability by item family and criterion conflict. A third should identify cases where the gold label shifted after adjudication, because those are likely taxonomy-instability cases rather than mere model failures.

A fourth table should report application-surface transfer. It should separate prompt-only items from webpage, document, email, tool-result, and transcript wrappers, with risk types split into integrity, confidentiality, tool misuse, policy bypass, and evaluator deception. This prevents colour-token success from being mistaken for agent-security robustness.

\section{Metrics}

The project starts with five metrics.

\subsection{Instruction-source sensitivity}

Instruction-source sensitivity measures how behaviour changes when the same string appears as user instruction, system-like text, webpage text, quoted passage, tool output, or transcript material.

\subsection{Mention/use robustness}

Mention/use robustness measures whether the model distinguishes discussing a string from enacting it.

\subsection{Scope robustness}

Scope robustness measures whether behaviour changes appropriately under negation, conditionals, modals, and exception clauses.

\subsection{Rater-role divergence}

Rater-role divergence measures systematic disagreement among policy evaluators, ordinary users, linguistic experts, domain specialists, and LLM judges.

\subsection{Taxonomy drift and diagnostic ambiguity}

Taxonomy drift measures whether labels shift under paraphrase, model version, item embedding, or evaluator type. Diagnostic ambiguity marks cases where disagreement is stable enough to be a property of the case, not just rater error.

\section{Implications for Safety Evaluation}

The benchmark is designed for practical diagnostic work. Model-policy teams can use it to find vague boundaries before those boundaries become grading failures. Red-teamers can use it to separate prompt-injection susceptibility from ordinary task failure. System-card authors can use it to expose fragile categories hidden by aggregate pass rates. External evaluators can use its provenance and claim register to ask exactly what a result licenses. These uses remain item-, procedure-, and configuration-specific unless their own validation and projection requirements are met (Section~\ref{sec:limitations}).

On the delegation-assurance account \parencite{reynolds2026delegationAssurance}, the benchmark supplies one measurement layer: evidence about whether model behaviour responds appropriately to instruction status, authority, scope, and failure attribution under controlled linguistic contrasts. It doesn't by itself establish institutional authorization, the adequacy of the resulting action record, or transport beyond the declared items, models, and wrappers; those require separate analyses. The companion evidentiary-assurance account \parencite{reynolds2026evidentiaryAssurance} treats the adequacy of that action record as its own layer. The three accounts divide the labour: this benchmark is the measurement layer, delegation assurance the normative-authorization layer, and evidentiary assurance the record-warrant layer.

The larger methodological point is that language-mediated distinctions are part of the control surface in these AI safety evaluations. A procedure that can't separate quotation from command, authority from content, policy analysis from unsafe enactment, or refusal from inability doesn't support the intended interpretation. A procedure that can make those separations on a fixed sample still needs separate evidence before its references, behaviours, measurements, or taxonomy are projected elsewhere.

\section{Conclusion}

Adversarial pragmatics reframes a common safety-evaluation problem. Many failures that look like model reasoning failures or policy failures are also failures of language-mediated control. The proposed benchmark, taxonomy, rater protocol, and drift metrics make those failures inspectable.

The project succeeds if it makes safety claims easier to audit: which string had authority, which policy boundary applied, which rater role produced the label, which transcript evidence supports the failure attribution, and which category becomes unstable under paraphrase or evaluator change.


% Restore a non-blind acknowledgements / AI-use statement only after the actual
% models used on the project have been recorded and current names verified.

\clearpage
\printbibliography

\end{document}
